\documentclass[11pt,letterpaper]{article}
\usepackage[T1]{fontenc}
\usepackage{lmodern}
\usepackage{graphicx}
\graphicspath{{./}{../}}
\providecommand{\pandocbounded}[1]{#1}
\usepackage{amsmath,amssymb,mathtools}
\usepackage[margin=1in,top=.85in,bottom=.85in]{geometry}
\usepackage{microtype}
\usepackage{enumitem}
\usepackage{titlesec}
\usepackage{needspace}
\usepackage[hidelinks]{hyperref}
\setlength{\parindent}{0pt}
\setlength{\parskip}{4.5pt plus 1pt minus 1pt}
\setlist{nosep,leftmargin=1.6em,topsep=4pt}
\setlength{\emergencystretch}{1.5em}
\allowdisplaybreaks[1]
\raggedbottom
\widowpenalty=10000
\clubpenalty=10000
\titleformat{\section}{\Large\bfseries}{}{0pt}{}
\titleformat{\subsubsection}{\normalsize\bfseries}{}{0pt}{}
\titleformat{\subsection}{\large\bfseries}{}{0pt}{}
\titlespacing*{\section}{0pt}{14pt}{8pt}
\titlespacing*{\subsubsection}{0pt}{9pt}{4pt}
\titlespacing*{\subsection}{0pt}{12pt}{6pt}
\newcommand{\tightlist}{\setlength{\itemsep}{0pt}\setlength{\parskip}{0pt}}
\hypersetup{pdftitle={Honours Linear Algebra - Weeks 1-3
Review},pdfsubject={Honours Linear Algebra: Week 4 / No lecture}}
\begin{document}
\begin{center}
{\LARGE Honours Linear Algebra - Weeks 1-3 Review\par}
\vspace{5pt}
{\large Week 4 / No lecture\par}
\end{center}
\vspace{4pt}
\Needspace{7\baselineskip}
\subsection{The same problem, three
languages}\label{the-same-problem-three-languages}

\(Ax=b\): equations / a combination of columns / a linear map
\(T(x)=Ax\).

\textbf{Basis = independent + spanning.} Existence of coordinates comes
from spanning; uniqueness from independence. For \(\dim V=n\), exactly
\(n\) vectors: either property implies both. Count matters.

\textbf{Rank--nullity.} \(A\in M_{m\times n}(K)\), rank \(r\): \[
\dim\ker T=n-r,\qquad
T\text{ injective}\iff r=n,\qquad T\text{ onto}\iff r=m.
\] For consistent \(Ax=b\), all solutions are \(x_p+\ker T\). A subspace
only when \(b=0\).

\textbf{Row reduction.} \(R=EA\), \(E\) invertible: same kernel, row
space and column relations. For a column-space basis, use pivot columns
of the \textbf{original} \(A\).

\Needspace{7\baselineskip}
\subsection{One matrix ties it
together}\label{one-matrix-ties-it-together}
\[
A=\begin{pmatrix}1&2&1\\0&0&1\end{pmatrix}
\sim R=\begin{pmatrix}1&2&0\\0&0&1\end{pmatrix}.
\] Pivots \(1,3\); \(a_2=2a_1\). Thus \[
\ker A=\operatorname{Span}\left\{\begin{pmatrix}-2\\1\\0\end{pmatrix}\right\},\qquad
\operatorname{Col}(A)=\operatorname{Span}\left\{\begin{pmatrix}1\\0\end{pmatrix},
\begin{pmatrix}1\\1\end{pmatrix}\right\}=\mathbb R^2.
\] \(T:\mathbb R^3\to\mathbb R^2\) is onto, not injective. Check
\(3=1+2\). \[
Ax=\begin{pmatrix}3\\1\end{pmatrix}
\iff x=\begin{pmatrix}2\\0\\1\end{pmatrix}
+t\begin{pmatrix}-2\\1\\0\end{pmatrix},\quad t\in\mathbb R.
\] The two columns in the image are basis vectors; the null vector is a
relation among all three columns. Different roles.

\Needspace{7\baselineskip}
\subsection{Direct sums; maps in
coordinates}\label{direct-sums-maps-in-coordinates}

\(V=U\oplus W\): every \(v=u+w\) uniquely. Subtract two decompositions:
\[
u-u'=w'-w\in U\cap W=\{0\}.
\qquad \dim(U+W)=\dim U+\dim W-\dim(U\cap W).
\] For three summands, pairwise zero intersections do \textbf{not}
suffice: the \(x\)-axis, \(y\)-axis and \(y=x\) in \(\mathbb R^2\).

\textbf{Matrix columns are basis images:} \[
[T]^\mathcal B_\mathcal C
=\bigl([T(b_1)]_\mathcal C\ \cdots\ [T(b_n)]_\mathcal C\bigr),\qquad
[T(v)]_\mathcal C=[T]^\mathcal B_\mathcal C[v]_\mathcal B.
\] Composition: rightmost acts first. The intermediate coordinate basis
must match.

\textbf{Keep the hypotheses.} Equal finite dimensions: injective
\(\iff\) onto. On \(\mathbb R[t]\), differentiation is onto but kills
constants. Also, \(\dim_{\mathbb R}\mathbb C=2\) while
\(\dim_{\mathbb C}\mathbb C=1\): always fix the field.
\end{document}
